% English translation of the verified Chinese baseline. Structural identifiers must remain synchronized.
\section{Symplectic Discretization, Birkhoff Structures, and Invariant-Torus Validation}\label{sec:symplectic-kam}

\subsection{Analysis branches}\label{sec:symplectic-kam-5-1}

A fixed-dimensional Galerkin truncation has already been completed before the KAM analysis. The principal theorem should therefore use a finite-dimensional Hamiltonian-flow theorem or a finite-dimensional symplectic-map theorem. Kuksin's infinite-dimensional KAM theory\cite{Kuksin1993InfiniteDimensionalKAM} is not used as a direct theorem for the fixed truncated system; it is retained only as a theoretical extension for discussions of the limit of infinitely many modes, the PDE limit, or consistency of the truncation. Moser's invariant-curve theorem for area-preserving annulus maps\cite{Moser1962InvariantCurves} primarily concerns two-dimensional twist maps. It may be used for low-dimensional sections and methodological comparison, but it cannot replace the Pöschel/Shang conditions required for multi-degree-of-freedom continuous Hamiltonian systems or high-dimensional symplectic maps.\cite{Shang1999KAMAlgorithms,Shang2000KAMMappings,Poschel2001KAMLecture}

Rather than describing every dynamical branch as a \(K\left( \mathbf{z},t \right)\)-symplectic system, this section distinguishes among:

\begin{itemize}
\tightlist
\item
  standard \(\mathbf{J}\)-symplectic discretization of an autonomous conservative Hamiltonian system;
\item
  standard nonautonomous or extended-phase-space symplectic discretization of an exact nonautonomous Hamiltonian system;
\item
  time-dependent \(\mathbf{K}\left( \mathbf{z},t \right)\)-symplectic discretization of a Birkhoff system that passes the variational self-adjointness test.
\end{itemize}

The Hamiltonian Birkhoff normal form is used in Sections 3 and 4 to construct \(H_{0} + \varepsilon H_{1}\). The self-adjoint Birkhoff representation considered here is used to test the structure of a general first-order system. The names are similar, but the mathematical objects are different.

\subsection{Standard symplectic branch and generating functions}\label{sec:symplectic-kam-5-2}

For the autonomous conservative CHNN

\begin{equation}
\label{eq:symplectic-kam-5-1}
\dot{\mathbf{Z}} = \mathbf{J}\nabla{\widehat{H}}_{\lambda}\left( \mathbf{Z} \right),
\end{equation}

the exact flow satisfies

\begin{equation}
\label{eq:symplectic-kam-5-2}
\left( D\Phi_{t} \right)^{T}\mathbf{J}\left( D\Phi_{t} \right) = \mathbf{J}.
\end{equation}

If the one-step map \(F_{h}^{cons}\) is an exact symplectic map sufficiently close to the identity and the relevant Jacobian block of the chosen projection is nondegenerate, a local generating function exists. Retaining the leading term of a center-type generating function gives the implicit midpoint rule

\begin{equation}
\label{eq:symplectic-kam-5-3}
\frac{\mathbf{Z}_{n + 1} - \mathbf{Z}_{n}}{h} = \mathbf{J}\nabla{\widehat{H}}_{\lambda}\left( \frac{\mathbf{Z}_{n + 1} + \mathbf{Z}_{n}}{2} \right),
\end{equation}

which is a second-order symplectic method. Higher-order methods may be constructed through generating-function expansions, composition methods, or other proven symplectic algorithms\cite{FengQin2010}.

For a nonautonomous Hamiltonian \(H\left( \mathbf{Z},t \right)\), introduce the extended phase space

\begin{equation}
\label{eq:symplectic-kam-5-4}
\mathcal{H}\left( \mathbf{Z},t,p_{t} \right) = H\left( \mathbf{Z},t \right) + p_{t},
\end{equation}

and integrate with respect to the extended canonical symplectic form. A transition to the formulation of the next subsection is needed only when the variables are noncanonical or when the system possesses a special Birkhoff representation.

The order-2 BDF scheme used for canonical data processing is not a symplectic stepper of the type defined in this section. It supplies only an observable \texttt{tsteps} sequence for COMSOL solver-step state sampling and canonical vector-field label generation. The frozen single-sentinel record can at most be a candidate input to a later A1 qualification experiment; it does not establish G0--G6 qualification across modes, amplitudes, or operating conditions. Neither BDF energy drift nor second-order convergence is used to infer symplecticity; a subsequent KAM computation must use a separately proven symplectic map and quantify its discretization defect.

\subsection{Self-adjoint representation of a Birkhoff system}\label{sec:symplectic-kam-5-3}

Let \(\mathbf{z} \in \mathbb{R}^{2n}\), let \(F_{i}\left( \mathbf{z},t \right)\) be Birkhoff functions, let \(B\left( \mathbf{z},t \right)\) be the Birkhoffian, and define

\begin{equation}
\label{eq:symplectic-kam-5-5}
K_{ij}\left( \mathbf{z},t \right) = \frac{\partial F_{j}}{\partial z_{i}} - \frac{\partial F_{i}}{\partial z_{j}}.
\end{equation}

The Birkhoff equations are

\begin{equation}
\label{eq:symplectic-kam-5-6}
\boxed{K_{ij}\left( \mathbf{z},t \right){\dot{z}}_{j} - \frac{\partial B}{\partial z_{i}} - \frac{\partial F_{i}}{\partial t} = 0.}
\end{equation}

Consider a general system written as

\begin{equation}
\label{eq:symplectic-kam-5-7}
K_{ij}{\dot{z}}_{j} + D_{i} = 0.
\end{equation}

On a local star-shaped domain, if \(\mathbf{K}\) is nondegenerate and satisfies skew symmetry, the spatial closure condition, and the temporal compatibility condition, then the system admits a regular self-adjoint Birkhoff representation. Equivalently, functions \(F_{i}\) and \(B\) can be found such that the preceding equation holds\cite{FengQin2010,HairerLubichWanner2006GNI}. Define the time-dependent two-form

For clarity, the closure and compatibility conditions used here are
\[
\partial_{z_k}K_{ij}+\partial_{z_i}K_{jk}+\partial_{z_j}K_{ki}=0,\qquad
\partial_t K_{ij}=\partial_{z_j}D_i-\partial_{z_i}D_j,
\]
on the same local domain, together with \(K_{ij}=-K_{ji}\) and \(\inf_{\mathcal{D}}\sigma_{min}(K)>0\). The present manuscript does not instantiate these conditions numerically.

\begin{equation}
\label{eq:symplectic-kam-5-8}
\Omega_{K} = \frac{1}{2}K_{ij}\left( \mathbf{z},t \right)\, dz_{i} \land dz_{j}.
\end{equation}

Its phase flow satisfies the time-dependent structure-preservation relation

\begin{equation}
\label{eq:symplectic-kam-5-9}
\left( D\Phi_{t,t_{0}} \right)^{T}\mathbf{K}\left( \Phi_{t,t_{0}},t \right)D\Phi_{t,t_{0}} = \mathbf{K}\left( \mathbf{z}_{0},t_{0} \right).
\end{equation}

When \(\mathbf{K} = \mathbf{J}^{- 1}\) is constant, this condition reduces to the canonical symplectic condition.

\subsection{Identification of a Birkhoff structure for the residual vector field}\label{sec:symplectic-kam-5-4}

The analysis is applied only to \(X_{res}\), as defined in \Cref{sec:chnn-model}. Parameterize the one-form directly as

\begin{equation}
\label{eq:symplectic-kam-5-10}
\mathcal{F}_{\rho} = \sum_{i = 1}^{2n}F_{\rho,i}\left( \mathbf{z},t \right)\, dz_{i},\quad\quad K_{\rho,ij} = \partial_{z_{i}}F_{\rho,j} - \partial_{z_{j}}F_{\rho,i},
\end{equation}

introduce \(B_{\beta}\), and define the candidate vector field

\begin{equation}
\label{eq:symplectic-kam-5-11}
{\widetilde{X}}_{\rho,\beta} = \mathbf{K}_{\rho}^{- 1}\left( \nabla B_{\beta} + \partial_{t}\mathbf{F}_{\rho} \right).
\end{equation}

The training loss is

\begin{equation}
\label{eq:symplectic-kam-5-12}
\mathcal{L}_{Birk} = \frac{1}{N}\sum_{\ell = 1}^{N}\left. \  \parallel \mathbf{K}_{\rho}X_{res} - \nabla B_{\beta} - \partial_{t}\mathbf{F}_{\rho} \right.\  \parallel_{2}^{2} + \mu_{nd}\frac{1}{N}\sum_{\ell = 1}^{N}\left\lbrack \max\{ 0,\kappa - \sigma_{min}\left( \mathbf{K}_{\rho} \right)\} \right\rbrack^{2}.
\end{equation}

On the domain \(\mathcal{D}\), define the defect of the self-adjoint Birkhoff representation by

\begin{equation}
\label{eq:symplectic-kam-5-13}
\boxed{\delta_{Birk} = \inf_{F,B,\ \inf_{\mathcal{D}}\sigma_{min}\left( K_{F} \right) \geq \kappa}\left. \  \parallel K_{F}X_{res} - \nabla B - \partial_{t}F \right.\  \parallel_{\mathcal{X}\left( \mathcal{D} \right)}.}
\end{equation}

If \(\delta_{Birk} = 0\), an exact representation is obtained. If \(\delta_{Birk} > 0\), only the projected vector field \(\widetilde{X}\) has an exact \(K\)-symplectic structure, and the remaining error must be controlled separately. The complete open system containing a nonstructured residual cannot therefore be described as a strictly \(K\left( \mathbf{z},t \right)\)-symplectic system.

\subsection{\texorpdfstring{A \(K\left( \mathbf{z},t \right)\)-symplectic discretization for Birkhoff systems}{A K-symplectic discretization for Birkhoff systems}}\label{sec:symplectic-kam-5-5}

Under regularity, self-adjointness, analyticity, and nondegeneracy conditions, the Birkhoff phase flow admits a gradient map and a scalar generating function under an appropriate local projection. If the generating function is analytic in the step size, it can be expanded as

\begin{equation}
\label{eq:symplectic-kam-5-14}
S\left( \mathbf{Z},h \right) = \sum_{r = 1}^{\infty}S_{r}\left( \mathbf{Z} \right)h^{r}.
\end{equation}

Truncation at order \(m\) gives an implicit order-\(m\) approximately \(K\left( \mathbf{z},t \right)\)-symplectic method satisfying

\begin{equation}
\label{eq:symplectic-kam-5-15}
\left( D\Psi_{h} \right)^{T}\mathbf{K}\left( \Psi_{h}\left( \mathbf{z} \right),t + h \right)D\Psi_{h} = \mathbf{K}\left( \mathbf{z},t \right) + O\left( h^{m + 1} \right).
\end{equation}

This result guarantees only approximate structure preservation and the stated order when its hypotheses hold; it is exactly \(K\)-symplectic only when the remainder vanishes on the domain of interest, and it does not by itself imply persistence of classical Hamiltonian KAM tori. If a nonstructured residual is superposed, the resulting structural defect must enter the error budget.

\subsection{Invariant-torus computation by the parameterization method}\label{sec:symplectic-kam-5-6}

For the exact symplectic one-step map \(F_{h}^{cons}\) of the autonomous conservative branch, seek an embedding

\begin{equation}
\label{eq:symplectic-kam-5-16}
K:\mathbb{T}^{d} \rightarrow \mathcal{M}
\end{equation}

satisfying

\begin{equation}
\label{eq:symplectic-kam-5-17}
\boxed{F_{h}^{cons}\left( K\left( \bm{\theta} \right) \right) = K\left( \bm{\theta} + \bm{\omega} \right).}
\end{equation}

Use the Fourier truncation

\begin{equation}
\label{eq:symplectic-kam-5-18}
K_{N}\left( \bm{\theta} \right) = \sum_{|k| \leq N}^{}{\widehat{K}}_{k}e^{2\pi ik \cdot \bm{\theta}},\quad\quad{\widehat{K}}_{- k} = \overline{{\widehat{K}}_{k}},
\end{equation}

and define the invariance residual

\begin{equation}
\label{eq:symplectic-kam-5-19}
E_{N}\left( \bm{\theta} \right) = F_{h}^{cons}\left( K_{N}\left( \bm{\theta} \right) \right) - K_{N}\left( \bm{\theta} + \bm{\omega} \right).
\end{equation}

FFT analysis, Poincaré sections, and long trajectories are used only to initialize \(K_{N}\) and \(\bm{\omega}\); a Newton-like method is then used for correction. The a posteriori KAM theorems of Haro et al. require the map to be analytic and exact symplectic, the frequency to satisfy a Diophantine condition, the dynamics to be nondegenerate, and \(\parallel E_{N} \parallel\) to satisfy an explicit smallness condition. Existence of a true invariant torus can be certified only after all constants have been verified. A small residual alone supports only the statement that a numerical approximate torus has been obtained\cite{Haro2016Parameterization}.

\subsection{Parameter continuation and safe operating domain}\label{sec:symplectic-kam-5-7}

Continuation is performed with respect to magnetic-field strength, insertion speed, flow velocity, contact clearance, and material parameters. At each parameter point, record:

\begin{itemize}
\tightlist
\item
  \(\varepsilon_{map}\) and the Hamiltonian learning error;
\item
  the symplectic or \(K\)-symplectic structural defect;
\item
  the invariance residual \(\parallel E_{N} \parallel\);
\item
  the minimum singular value of the twist matrix/frequency map;
\item
  the finite-order Diophantine margin and decay of the Fourier coefficients;
\item
  the energy balance, port power, peak contact response, and trajectory error;
\item
  the discrepancy between the Kirchhoff model and the high-fidelity model.
\end{itemize}

Define the KAM-theoretic guarantee domain as

\begin{equation}
\label{eq:symplectic-kam-5-20}
\mathcal{S}_{KAM} = \{\mu:\text{\footnotesize the symplecticity, nondegeneracy, nonresonance, and a posteriori smallness conditions all hold}\}.
\end{equation}

The complete engineering safe domain is defined as

\begin{equation}
\label{eq:symplectic-kam-5-21}
\boxed{\mathcal{S}_{safe} = \mathcal{S}_{KAM} \cap \mathcal{S}_{energy} \cap \mathcal{S}_{contact} \cap \mathcal{S}_{model} \cap \mathcal{S}_{control}.}
\end{equation}

Here, KAM theory provides only a conditional long-time stability guarantee for the conservative backbone. Open ports, dissipation, contact switching, and high-fidelity model discrepancy are addressed by additional indicators. The safe domain is neither a single ``chaos threshold'' nor evidence of clinical efficacy.

At present, \(\mathcal{S}_{KAM}\), \(\mathcal{S}_{safe}\), the invariant-torus embedding \(K_N\), and its a posteriori constants have not been instantiated numerically. The complete algorithms and criteria in this section define the subsequent numerical validation program; the current results do not constitute numerical KAM certification and do not establish a true invariant torus or an engineering safe domain.
